% ৯. Indexing
\section{Indexing}
ইনডেক্সিং হলো বইয়ের সূচিপত্রের মতো একটি সহায়ক ব্যবস্থা, যা ডেটাবেজে নির্দিষ্ট রেকর্ড দ্রুত খুঁজে পেতে সাহায্য করে। বড় টেবিলে বারবার অনুসন্ধান করা হলে ইনডেক্স খুব কার্যকর।

\begin{definitionbox}{Indexing}
কোনো টেবিলের নির্দিষ্ট ফিল্ডের মানের ভিত্তিতে দ্রুত অনুসন্ধানের জন্য তৈরি বিশেষ ডেটা কাঠামোকে Index বলে; এই প্রক্রিয়াকে Indexing বলে।
\end{definitionbox}

\begin{keypointbox}{সুবিধা ও সীমাবদ্ধতা}
\begin{itemize}
    \item সুবিধা: অনুসন্ধান দ্রুত হয়, Sort ও Join অপারেশনে সহায়তা করে।
    \item সীমাবদ্ধতা: অতিরিক্ত স্টোরেজ লাগে এবং Insert/Update/Delete কিছুটা ধীর হতে পারে।
\end{itemize}
\end{keypointbox}

\begin{examplebox}{ব্যবহার}
StudentID, Roll, NID, AccountNo-এর মতো অনন্য বা বেশি অনুসন্ধানযোগ্য ফিল্ডে ইনডেক্স রাখা উপযোগী। কিন্তু Address বা Comment-এর মতো বড় লেখা ফিল্ডে সাধারণত ইনডেক্স কম দরকার।
\end{examplebox}

\begin{mistakebox}{Sorting বনাম Indexing}
Sorting ডেটা প্রদর্শনের ক্রম বদলায়। Indexing মূলত অনুসন্ধান দ্রুত করার জন্য সহায়ক কাঠামো তৈরি করে। দুইটি একই বিষয় নয়।
\end{mistakebox}
